% Chapter 3: System Analysis
\chapter{SYSTEM ANALYSIS}

\section{Introduction}

System analysis is a critical phase in software development that involves understanding the problem domain, gathering requirements, and determining feasibility. This chapter presents the requirements analysis, feasibility study, and system specifications for EduSphere AI.

\section{Requirements Analysis}

\subsection{Functional Requirements}

Functional requirements define what the system should do. The following functional requirements are identified for EduSphere AI:

\subsubsection{User Management}
\begin{enumerate}[label=FR-UM-\arabic*]
\item System shall support user registration with email and password
\item System shall assign one of six roles: Student, Faculty, HOD, COE, Principal, Admin
\item System shall require HOD approval for new user creation
\item System shall require Principal approval after HOD approval
\item System shall allow profile updates by users
\item System shall support password reset functionality
\item System shall maintain active/suspended/rejected user status
\end{enumerate}

\subsubsection{Document Management}
\begin{enumerate}[label=FR-DM-\arabic*]
\item System shall allow document upload in PDF and Office formats
\item System shall capture document metadata (title, type, category, scope, department)
\item System shall store documents securely in cloud storage
\item System shall generate unique document identifiers
\item System shall support document search by title, type, and date
\item System shall maintain document version history
\item System shall allow document creators to view their submissions
\end{enumerate}

\subsubsection{Workflow Management}
\begin{enumerate}[label=FR-WF-\arabic*]
\item System shall analyze document context (type, scope, department)
\item System shall generate approval workflows automatically
\item System shall create approval steps in sequence
\item System shall support workflows with 1-5 approval steps
\item System shall match approvers by role, not by specific user ID
\item System shall track workflow status (pending, approved, rejected)
\item System shall auto-advance workflows upon approval
\item System shall halt workflows upon rejection
\end{enumerate}

\subsubsection{Digital Signatures}
\begin{enumerate}[label=FR-SIG-\arabic*]
\item System shall provide signature capture interface (draw/upload)
\item System shall allow signature placement on PDF documents
\item System shall embed signatures directly into PDF files
\item System shall store signatures securely with encryption
\item System shall timestamp all signatures
\item System shall link signatures to specific approval steps
\item System shall prevent signature forgery through validation
\end{enumerate}

\subsubsection{Assignment Management}
\begin{enumerate}[label=FR-ASN-\arabic*]
\item System shall support four assignment types: code, link, project, document
\item System shall allow faculty to create assignments with deadlines
\item System shall enable students to submit assignments before deadline
\item System shall execute code submissions with test cases
\item System shall auto-grade code assignments based on test results
\item System shall track link assignment access
\item System shall support team assignments with multiple members
\item System shall allow faculty to manually grade submissions
\item System shall calculate and display leaderboards
\end{enumerate}

\subsubsection{Notice Board}
\begin{enumerate}[label=FR-NOT-\arabic*]
\item System shall allow authorized users to create notices
\item System shall categorize notices (Academic, Events, Holidays, etc.)
\item System shall support targeted publishing by role/department
\item System shall attach files to notices
\item System shall require approval before publication
\item System shall display published notices on notice board
\item System shall allow search and filtering of notices
\end{enumerate}

\subsubsection{Department Management}
\begin{enumerate}[label=FR-DEPT-\arabic*]
\item System shall allow Principal/Admin to create departments
\item System shall allow creation of sections within departments
\item System shall assign users to departments and sections
\item System shall filter data based on user's department
\item System shall support department-wise analytics
\end{enumerate}

\subsubsection{Admissions Management}
\begin{enumerate}[label=FR-ADM-\arabic*]
\item System shall manage admission seasons
\item System shall support semester promotion workflows
\item System shall handle fresh admission registrations
\item System shall track fee payments and installments
\item System shall maintain student admission history
\item System shall require multiple approval steps
\end{enumerate}

\subsection{Non-Functional Requirements}

Non-functional requirements specify how the system should perform.

\subsubsection{Performance}
\begin{enumerate}[label=NFR-PERF-\arabic*]
\item Page load time shall be under 2 seconds for 95\% of requests
\item System shall support 500 concurrent users
\item Database queries shall execute in under 500ms
\item Code execution results shall return within 10 seconds
\item File uploads shall support files up to 25MB
\end{enumerate}

\subsubsection{Security}
\begin{enumerate}[label=NFR-SEC-\arabic*]
\item System shall use HTTPS for all communications
\item System shall hash passwords using bcrypt with salt
\item System shall implement JWT-based authentication
\item System shall enforce row-level security in database
\item System shall validate all user inputs
\item System shall prevent SQL injection attacks
\item System shall protect against XSS attacks
\item System shall implement CSRF protection
\end{enumerate}

\subsubsection{Scalability}
\begin{enumerate}[label=NFR-SCAL-\arabic*]
\item System shall scale horizontally to handle increased load
\item Database shall support read replicas
\item File storage shall use CDN for global distribution
\item System architecture shall be cloud-native
\end{enumerate}

\subsubsection{Usability}
\begin{enumerate}[label=NFR-USE-\arabic*]
\item Interface shall be intuitive requiring minimal training
\item System shall be responsive on mobile devices
\item Error messages shall be clear and actionable
\item System shall provide inline help and tooltips
\item Color schemes shall follow accessibility standards (WCAG 2.1)
\end{enumerate}

\subsubsection{Reliability}
\begin{enumerate}[label=NFR-REL-\arabic*]
\item System uptime shall be 99.5\% or higher
\item Database shall have automated daily backups
\item System shall recover gracefully from failures
\item Data shall be replicated across multiple zones
\end{enumerate}

\subsubsection{Maintainability}
\begin{enumerate}[label=NFR-MAIN-\arabic*]
\item Code shall follow TypeScript strict mode
\item System shall use consistent naming conventions
\item Code shall be modular and component-based
\item System shall log errors for debugging
\item Documentation shall be comprehensive
\end{enumerate}

\section{Feasibility Study}

\subsection{Technical Feasibility}

\textbf{Hardware Requirements:}
\begin{itemize}
\item Development: Standard laptop/desktop with 8GB RAM, modern processor
\item Deployment: Cloud infrastructure (Vercel, Supabase)
\item Client Access: Any device with modern web browser
\end{itemize}

\textbf{Software Requirements:}
\begin{itemize}
\item Node.js 18+ for Next.js runtime
\item PostgreSQL database (provided by Supabase)
\item Git for version control
\item VS Code or similar IDE for development
\end{itemize}

\textbf{Technical Expertise:}
\begin{itemize}
\item Web development with JavaScript/TypeScript
\item React and Next.js framework knowledge
\item Database design and SQL
\item RESTful API development
\item Cloud platform familiarity
\end{itemize}

\textbf{Conclusion:} The project is technically feasible. All required technologies are mature, well-documented, and accessible. The development team possesses necessary skills, and learning resources are abundantly available.

\subsection{Economic Feasibility}

\textbf{Development Costs:}
\begin{itemize}
\item Development tools: \$0 (using free/open-source software)
\item Hosting during development: \$0 (Supabase free tier, Vercel hobby plan)
\item Domain name: \$10-15/year
\item SSL certificate: \$0 (Let's Encrypt)
\end{itemize}

\textbf{Operational Costs (estimated for 1000 users):}
\begin{itemize}
\item Supabase Pro: \$25/month
\item Vercel Pro: \$20/month
\item Judge0 API: \$0-50/month depending on usage
\item Total: \$45-95/month (\$540-1140/year)
\end{itemize}

\textbf{Comparison with Commercial Solutions:}
\begin{itemize}
\item SharePoint: \$5-15 per user/month = \$5000-15000/month for 1000 users
\item DocuSign: \$10+ per user/month
\item Custom Development: \$50,000-100,000 upfront
\end{itemize}

\textbf{Return on Investment:}
\begin{itemize}
\item Time savings: 10-20 hours/week for administrative staff
\item Paper reduction: \$1000+/year
\item Improved efficiency: Faster approvals, better accountability
\item Payback period: Immediate (development costs are minimal)
\end{itemize}

\textbf{Conclusion:} The project is highly economically feasible. Development and operational costs are minimal compared to commercial alternatives while providing comprehensive functionality.

\subsection{Operational Feasibility}

\textbf{User Acceptance:}
\begin{itemize}
\item System reduces manual workload
\item Provides transparency and real-time tracking
\item Accessible from any device
\item Modern, intuitive interface
\item Reduces physical presence requirements
\end{itemize}

\textbf{Organizational Impact:}
\begin{itemize}
\item Minimal disruption: Gradual rollout possible
\item Training requirements: Low (intuitive interface)
\item Process changes: Mirrors existing workflows
\item Support requirements: Basic IT support sufficient
\end{itemize}

\textbf{Legal and Compliance:}
\begin{itemize}
\item Digital signatures: Legally valid under IT Act 2000 (India)
\item Data privacy: Compliant with institutional policies
\item Audit trails: Exceed manual system capabilities
\item Data ownership: Institution retains full control
\end{itemize}

\textbf{Conclusion:} The system is operationally feasible. User acceptance is expected to be high due to clear benefits. Legal requirements are met, and organizational disruption is minimal.

\subsection{Schedule Feasibility}

\textbf{Development Timeline:}
\begin{enumerate}
\item Requirements Analysis: 2 weeks
\item System Design: 2 weeks
\item Core Development: 8 weeks
\item Testing: 2 weeks
\item Deployment: 1 week
\item Total: 15 weeks (approximately 4 months)
\end{enumerate}

\textbf{Resource Availability:}
\begin{itemize}
\item Development team: Available full-time
\item Testing participants: Faculty and students willing to participate
\item Infrastructure: Available on-demand (cloud)
\end{itemize}

\textbf{Conclusion:} The project schedule is feasible within the academic year timeframe. The timeline allows for iterative development and adequate testing.

\section{System Specifications}

\subsection{Hardware Specifications}

\textbf{Server Requirements (Cloud Infrastructure):}
\begin{itemize}
\item CPU: 2+ vCPUs (auto-scaling)
\item RAM: 4GB+ (auto-scaling)
\item Storage: 50GB+ (expandable)
\item Network: 100Mbps+ bandwidth
\item Geographic distribution: Multi-region for redundancy
\end{itemize}

\textbf{Client Requirements:}
\begin{itemize}
\item Any device with modern web browser
\item Minimum 2GB RAM recommended
\item Internet connection: 1Mbps+ for smooth experience
\item Screen resolution: 360px+ width (mobile-responsive)
\end{itemize}

\subsection{Software Specifications}

\textbf{Development Environment:}
\begin{itemize}
\item Operating System: Windows 10+, macOS 10.15+, or Linux
\item Node.js: v18.0.0 or higher
\item npm: v9.0.0 or higher
\item Git: v2.30.0 or higher
\item IDE: VS Code with TypeScript, ESLint extensions
\end{itemize}

\textbf{Production Environment:}
\begin{itemize}
\item Runtime: Node.js 18+ (Vercel managed)
\item Database: PostgreSQL 15+ (Supabase managed)
\item Web Server: Next.js built-in server
\item File Storage: Supabase Storage (S3-compatible)
\item CDN: Vercel Edge Network
\end{itemize}

\textbf{Technology Stack:}
\begin{itemize}
\item Frontend: Next.js 14, React 18, TypeScript 5.7, Tailwind CSS 3.4
\item Backend: Next.js API Routes, Supabase
\item Database: PostgreSQL with Row-Level Security
\item Authentication: Supabase Auth (JWT-based)
\item File Storage: Supabase Storage
\item Code Execution: Judge0 API
\item PDF Manipulation: pdf-lib
\item Version Control: Git + GitHub
\end{itemize}

\section{Use Case Analysis}

\subsection{Primary Actors}
\begin{itemize}
\item Student
\item Faculty
\item Head of Department (HOD)
\item Chief of Examinations (COE)
\item Principal
\item System Administrator
\end{itemize}

\subsection{Key Use Cases}

\subsubsection{UC1: Upload Document}
\textbf{Actor:} Faculty\\
\textbf{Precondition:} User is authenticated and authorized\\
\textbf{Main Flow:}
\begin{enumerate}
\item User navigates to Upload Document page
\item User fills document metadata form
\item User selects file to upload
\item User submits form
\item System validates inputs
\item System analyzes document context
\item System generates approval workflow
\item System stores document and creates workflow records
\item System notifies first approver
\item System displays success message
\end{enumerate}
\textbf{Postcondition:} Document is uploaded and workflow initiated

\subsubsection{UC2: Approve Document}
\textbf{Actor:} HOD/Principal\\
\textbf{Precondition:} Document awaiting actor's approval\\
\textbf{Main Flow:}
\begin{enumerate}
\item User views pending approvals
\item User clicks on document to review
\item System displays document details and content
\item User captures/uploads digital signature
\item User provides optional comments
\item User clicks Approve button
\item System validates signature
\item System embeds signature (for PDF)
\item System updates approval status
\item System advances workflow
\item System notifies next approver (if any)
\item System displays confirmation
\end{enumerate}
\textbf{Postcondition:} Document approved and workflow advanced

\subsubsection{UC3: Submit Assignment}
\textbf{Actor:} Student\\
\textbf{Precondition:} Assignment exists and deadline not passed\\
\textbf{Main Flow:}
\begin{enumerate}
\item Student views assignment details
\item For code assignment:
    \begin{enumerate}
    \item Student writes/pastes code
    \item Student selects programming language
    \item Student clicks Run Code
    \item System executes code with test cases
    \item System displays results
    \item Student reviews and submits
    \end{enumerate}
\item For other types:
    \begin{enumerate}
    \item Student uploads file or provides link
    \item Student clicks Submit
    \end{enumerate}
\item System records submission
\item System calculates score (for auto-graded)
\item System updates leaderboard
\item System displays confirmation
\end{enumerate}
\textbf{Postcondition:} Assignment submitted and graded (if applicable)

\section{Conclusion}

The system analysis phase establishes clear requirements, confirms feasibility across all dimensions, and defines precise specifications. The functional and non-functional requirements provide a comprehensive blueprint for system design and implementation. Feasibility studies confirm that the project is viable from technical, economic, operational, and schedule perspectives. The specifications ensure all stakeholders have a common understanding of system capabilities and constraints.

\newpage
